% =====================================================================
% draft-product-trees-elbo.tex
%
% DRAFT (not part of the live published spec). A standalone derivation of
% the "product-of-trees" structured mean-field ELBO for the dynamic-latent-
% field (dynfield) coevolution cluster likelihood. This bound is valid for
% clusters of ANY size m and is LINEAR in m (never A^m), which is its whole
% reason for existing: the exact L*A^m field-augmented peel is intractable
% for m>2 at realistic A=20.
%
% This file deliberately does NOT \input into supplement.tex /
% appendix-tkfdp.tex (the live spec). It is a self-contained working
% derivation for the product-of-trees paper result. Compile standalone:
%     pdflatex draft-product-trees-elbo.tex
% =====================================================================
\documentclass[11pt]{article}
\usepackage{amsmath,amssymb,amsthm}
\usepackage[margin=1in]{geometry}
\usepackage{bm}
\usepackage{hyperref}

\newtheorem{proposition}{Proposition}
\newtheorem{lemma}{Lemma}
\newtheorem{corollary}{Corollary}
\newtheorem{remark}{Remark}

\newcommand{\btheta}{\bm{\theta}}
\newcommand{\bDelta}{\bm{\Delta}}
\newcommand{\bx}{\bm{x}}
\newcommand{\Ex}[1]{\mathbb{E}\!\left[#1\right]}
\newcommand{\Eq}[2]{\mathbb{E}_{#1}\!\left[#2\right]}
\newcommand{\KL}[2]{\mathrm{KL}\!\left(#1\,\|\,#2\right)}
\newcommand{\Ent}[1]{\mathcal{H}\!\left[#1\right]}
\newcommand{\pa}{\mathrm{pa}}
\newcommand{\Psub}{P_{\mathrm{sub}}}
\newcommand{\Real}{\mathbb{R}}

\title{The Product-of-Trees Structured Mean-Field ELBO\\for the Dynamic-Latent-Field Coevolution Model}
\author{tkf-dp working note}
\date{\today}

\begin{document}
\maketitle

\begin{abstract}
We derive a structured mean-field variational lower bound on the marginal
likelihood of a dynfield coevolution cluster of arbitrary size $m$. The
variational family is a \emph{product of independent tree-structured
factors}: one tree-Markov law over the field-and-jumps configuration, and,
for each of the $m$ sites, a tree-Markov law over that site's residue
trajectory with full along-tree correlation. Unlike the exact
field-augmented peel, whose per-node state space is $L\cdot A^{m}$, this
bound costs $O\!\big(m\,N\,A^{2}+N\,L^{2}\big)$ --- \emph{linear} in $m$.
We give the two block-coordinate-ascent updates in closed form (the residue
update uses a \emph{geometric-mean} branch operator, the field update folds
expected residue evidence into the field MRF potentials), prove the bound is
monotone under coordinate ascent, identify the regimes in which it is tight
or exact, and characterise the single approximation it makes: it severs the
posterior coupling between the field and the residues.
\end{abstract}

\section{Model and notation}

Fix a rooted tree with node set $\mathcal V$, root $r$, and for each
non-root node $v$ a parent $u=\pa(v)$ and branch length $\tau_v$. Let
$L=L_{\mathrm{field}}$ be the number of field states and $A$ the residue
alphabet size. A cluster of $m$ sites carries, at every node $v$,
\begin{itemize}
\item a \emph{field} state $\theta_v\in\{0,\dots,L-1\}$;
\item for each non-root $v$, a \emph{jump indicator} $\Delta_v\in\{0,1\}$ on
      the edge $u\to v$;
\item residue states $x_{n,v}\in\{0,\dots,A-1\}$, one per site $n=1,\dots,m$.
\end{itemize}
Write $\btheta=(\theta_v)_v$, $\bDelta=(\Delta_v)_{v\neq r}$, and
$\bx_n=(x_{n,v})_v$ for the whole residue trajectory of site $n$. Site $n$
has class $c_n$; the class maps a field state to an archetype
$a_n(\theta)$, whose residue stationary is $\pi(c_n,\theta,\cdot)$ and whose
GTR generator is $Q(c_n,\theta)$ with exchangeability $S$. Let
$\Psub^{c_n,\theta}(\tau)=\exp\!\big(Q(c_n,\theta)\tau\big)$.

The field is an F81-on-DP CTMC with rate $\rho_{\mathrm{chain}}$ and
stationary $\rho$. Per branch of length $\tau$ define, exactly as in the code
(\texttt{mm\_clv.field\_transition\_row}, \texttt{beta\_no\_jump}),
\begin{equation}
\begin{aligned}
P^{\Theta}_\tau(\theta_u,\theta_v)
  &= e^{-\rho_{\mathrm{chain}}\tau}\,[\theta_u{=}\theta_v]
   + (1-e^{-\rho_{\mathrm{chain}}\tau})\,\rho(\theta_v),\\
\beta_\tau(\theta) &= e^{-\rho_{\mathrm{chain}}(1-\rho(\theta))\tau},\qquad
W_\tau(\theta_u,\theta_v) = P^{\Theta}_\tau(\theta_u,\theta_v)
   - \beta_\tau(\theta_u)\,[\theta_u{=}\theta_v].
\end{aligned}
\label{eq:field-kernels}
\end{equation}

\paragraph{Jump-augmented joint.}
We take as the target the $\Delta$-augmented joint
\begin{equation}
P(\btheta,\bDelta,\bx)
  = \underbrace{\rho(\theta_r)\prod_{n}\pi(c_n,\theta_r,x_{n,r})}_{\text{root prior}}
    \;\prod_{v\neq r}\psi_v(\theta_u,\Delta_v,\theta_v,\bx_{\cdot,u},\bx_{\cdot,v}),
\label{eq:joint}
\end{equation}
with the branch factor
\begin{equation}
\psi_v =
\begin{cases}
\beta_{\tau_v}(\theta_u)\,[\theta_v{=}\theta_u]\;
   \prod_{n}\Psub^{c_n,\theta_u}(\tau_v)(x_{n,u},x_{n,v}), & \Delta_v=0\ \text{(no jump)},\\[4pt]
W_{\tau_v}(\theta_u,\theta_v)\;\prod_{n}\pi(c_n,\theta_v,x_{n,v}), & \Delta_v=1\ \text{(jump / renewal)}.
\end{cases}
\label{eq:psi}
\end{equation}
Summing \eqref{eq:psi} over $\Delta_v$ recovers the un-augmented branch factor
$K_v$ used by the exact peel
(\texttt{elbo\_rescorr.branch\_factor},
\texttt{exact\_cap2.branch\_pair}); hence
$\sum_{\bDelta}P(\btheta,\bDelta,\bx)=P(\btheta,\bx)$, and marginalising the
internal residues and all fields gives exactly the model likelihood
$P(\mathrm{obs})$ that the exact peel computes. Leaves clamp $x_{n,v}$ to the
observation (or leave it free at a gap).

\section{The product-of-trees variational family}

\begin{equation}
q(\btheta,\bDelta,\bx)
  = q_\Theta(\btheta,\bDelta)\;\prod_{n=1}^{m} q_n(\bx_n),
\label{eq:family}
\end{equation}
where
\begin{itemize}
\item $q_\Theta$ is an arbitrary \emph{tree-Markov} law over the
      field-and-jumps configuration --- exactly the tree MRF that
      \texttt{field\_bp.field\_logbp} handles, with a root node potential
      over $\theta_r$ and per-edge potentials over
      $(\theta_u,\theta_v,\Delta_v)$;
\item each $q_n$ is an arbitrary \emph{tree-Markov} law over site $n$'s
      residue trajectory (an $A$-state tree chain with full parent--child
      correlation), with leaves clamped to their observations.
\end{itemize}
The $m+1$ factors are \emph{independent}: residues are \emph{not} conditioned
on the field. This is the sole structural difference from the previous
per-site bound (\texttt{elbo\_persite}), whose family
$q_\Theta\prod_n q_n(\bx_n\mid\btheta,\bDelta)$ ties each residue chain to the
field and thereby pays $O(mLA^2)$ but keeps the field--residue seam. Dropping
that conditioning is what buys the clean product form below; Section
\ref{sec:tightness} quantifies what it costs.

Because $q$ factorises as in \eqref{eq:family}, its entropy is additive,
\begin{equation}
\Ent{q} = \Ent{q_\Theta}+\sum_{n=1}^m\Ent{q_n}.
\label{eq:entropy-split}
\end{equation}

\section{The ELBO}

For any $q$, Jensen gives the evidence lower bound
\begin{equation}
\log P(\mathrm{obs}) \;\ge\; \mathcal L(q)
   := \Eq{q}{\log P(\btheta,\bDelta,\bx)} + \Ent{q}
   \;=\; \log P(\mathrm{obs}) - \KL{q}{P(\cdot\mid\mathrm{obs})},
\label{eq:elbo-def}
\end{equation}
with equality iff $q$ is the exact posterior. Since $q$ in \eqref{eq:family}
is a proper normalised distribution and \eqref{eq:joint} normalises to
$P(\mathrm{obs})$, $\mathcal L(q)\le\log P(\mathrm{obs})$ \emph{always}: the
product-of-trees functional is a genuine lower bound. We now make every term
explicit in tree marginals.

Let $b_v(\theta)=q_\Theta(\theta_v{=}\theta)$ and
$e_v(\theta_u,\theta_v,\Delta)=q_\Theta(\theta_u,\theta_v,\Delta_v{=}\Delta)$
be the field node/edge marginals (returned by \texttt{field\_logbp}), and let
$\mu_{n,v}(x)=q_n(x_{n,v}{=}x)$, $\xi_{n,v}(x_u,x_v)=q_n(x_{n,u}{=}x_u,x_{n,v}{=}x_v)$
be the residue node/edge marginals. Using the factorisation
$q_\Theta\!\perp\! q_n$ and $q_n\!\perp\! q_{n'}$,
\begin{align}
\Eq{q}{\log P}
  =\;& \sum_\theta b_r(\theta)\Big[\log\rho(\theta)
        + \sum_n\sum_x \mu_{n,r}(x)\log\pi(c_n,\theta,x)\Big]
        \label{eq:ElogP-root}\\
  &+\sum_{v\neq r}\Bigg\{
     \sum_\theta e_v(\theta,\theta,0)\Big[\log\beta_{\tau_v}(\theta)
        + \sum_n\sum_{x_u,x_v}\xi_{n,v}(x_u,x_v)\log\Psub^{c_n,\theta}(\tau_v)(x_u,x_v)\Big]
        \nonumber\\
  &\hphantom{+\sum_{v\neq r}\Bigg\{}
   + \sum_{\theta_u,\theta_v} e_v(\theta_u,\theta_v,1)\Big[\log W_{\tau_v}(\theta_u,\theta_v)
        + \sum_n\sum_{x}\mu_{n,v}(x)\log\pi(c_n,\theta_v,x)\Big]\Bigg\}.
        \label{eq:ElogP-edge}
\end{align}
The $\Delta{=}0$ term survives only on the diagonal $\theta_u{=}\theta_v{=}\theta$
because $\psi_v$ forces $\theta_v{=}\theta_u$ when there is no jump.

Each residue entropy is the rooted tree-chain (junction-tree) entropy
\begin{equation}
\Ent{q_n} = \Ent{\mu_{n,r}} + \sum_{v\neq r}
  \Big(-\!\!\sum_{x_u,x_v}\!\xi_{n,v}(x_u,x_v)
        \big[\log\xi_{n,v}(x_u,x_v)-\log\mu_{n,u}(x_u)\big]\Big),
\label{eq:Hqn}
\end{equation}
which is exact for any tree-Markov $q_n$; at a leaf $v$ the residue is clamped,
so $\mu_{n,v}$ is a point mass and that edge contributes $0$.
$\Ent{q_\Theta}$ is delivered by \texttt{field\_logbp} through the identity in
Proposition~\ref{prop:field-elbo}.

\section{Block-coordinate-ascent updates}

We maximise $\mathcal L$ by alternating exact block updates.

\subsection{Residue update (fix $q_\Theta$, optimise each $q_n$)}
\label{sec:res-update}

Collect the $\bx_n$-dependent part of $\mathcal L$. From
\eqref{eq:ElogP-root}--\eqref{eq:ElogP-edge}, $\mathcal L$ is, up to a
$q_n$-independent constant,
\begin{equation}
\mathcal L = \Eq{q_n}{\;\log\phi_{n,r}(x_{n,r})
   + \sum_{v\neq r}\big[\,\log\Psi_{n,v}(x_{n,u},x_{n,v}) + \log\phi_{n,v}(x_{n,v})\,\big]}
   + \Ent{q_n} + \mathrm{const},
\label{eq:res-objective}
\end{equation}
with the tree pairwise-MRF log-potentials
\begin{align}
\log\phi_{n,r}(x) &= \sum_\theta b_r(\theta)\,\log\pi(c_n,\theta,x)
   &&\text{(root node potential)},\label{eq:phi-root}\\
\log\phi_{n,v}(x) &= \sum_{\theta_v} w^{1}_{v}(\theta_v)\,\log\pi(c_n,\theta_v,x),\quad
   w^{1}_v(\theta_v):=\!\sum_{\theta_u}\! e_v(\theta_u,\theta_v,1)
   &&\text{(node potential from renewal)},\label{eq:phi-node}\\
\log\Psi_{n,v}(x_u,x_v) &= \sum_\theta w^{0}_v(\theta)\,
   \log\Psub^{c_n,\theta}(\tau_v)(x_u,x_v),\quad
   w^{0}_v(\theta):=e_v(\theta,\theta,0)
   &&\text{(edge potential).}\label{eq:psi-edge}
\end{align}
\eqref{eq:res-objective} is $-\KL{q_n}{q_n^\star}+\mathrm{const}$ where
$q_n^\star\propto\phi_{n,r}\prod_v\Psi_{n,v}\phi_{n,v}$ is itself
tree-Markov; hence the \emph{exact} block maximiser is $q_n=q_n^\star$,
attained \emph{within} the family. Its marginals $\mu_{n,\cdot},\xi_{n,\cdot}$
and its entropy are obtained by one sum--product (forward--backward) sweep on
the site-$n$ tree, with leaves clamped to observations.

\begin{remark}[Geometric-mean branch operator]
The edge factor is
$\Psi_{n,v}(x_u,x_v)=\exp\!\big(\sum_\theta w^0_v(\theta)\log\Psub^{c_n,\theta}(\tau_v)(x_u,x_v)\big)
 = \prod_\theta\big[\Psub^{c_n,\theta}(\tau_v)\big]^{\,w^0_v(\theta)}(x_u,x_v)$,
a \emph{weighted geometric mean} (exp of the $q_\Theta$-weighted expected log)
of the per-field substitution operators --- \emph{not} the arithmetic mean.
This is the standard variational-message-passing form: the residue chain sees
each field state's transition matrix weighted \emph{multiplicatively} by the
joint no-jump probability $w^0_v(\theta)=q_\Theta(\theta_u{=}\theta_v{=}\theta,\Delta_v{=}0)$.
Using the arithmetic mean here would not be the coordinate maximiser and would
break the bound's monotonicity.
\end{remark}

\subsection{Field update (fix all $q_n$, optimise $q_\Theta$)}
\label{sec:field-update}

Collect the $(\btheta,\bDelta)$-dependent part of $\mathcal L$. It is the
energy of a field-and-jumps tree MRF with the potentials that
\emph{absorb the expected residue evidence}:
\begin{align}
\log\Phi_r(\theta) &= \log\rho(\theta) + \sum_n\sum_x\mu_{n,r}(x)\log\pi(c_n,\theta,x),
   \label{eq:field-root}\\
\log\Phi_v(\theta_u,\theta_v,0) &=
  \begin{cases}
   \log\beta_{\tau_v}(\theta) + \sum_n\sum_{x_u,x_v}\xi_{n,v}(x_u,x_v)\log\Psub^{c_n,\theta}(\tau_v)(x_u,x_v),
     & \theta_u{=}\theta_v{=}\theta,\\
   -\infty, & \text{otherwise},
  \end{cases}\label{eq:field-edge0}\\
\log\Phi_v(\theta_u,\theta_v,1) &=
   \log W_{\tau_v}(\theta_u,\theta_v) + \sum_n\sum_x\mu_{n,v}(x)\log\pi(c_n,\theta_v,x).
   \label{eq:field-edge1}
\end{align}
Running \texttt{field\_logbp} on these potentials returns the Gibbs law
$q_\Theta^\star\propto\exp(\text{energy})$, its marginals $b,e$, and
$\log Z_{\mathrm{field}}$. As in Section~\ref{sec:res-update}, $q_\Theta^\star$
is the exact block maximiser.

\begin{proposition}[Field partition function realises the field-side ELBO]
\label{prop:field-elbo}
With the potentials \eqref{eq:field-root}--\eqref{eq:field-edge1}, the
belief-propagation partition function satisfies
\begin{equation}
\log Z_{\mathrm{field}} = \Eq{q}{\log P} + \Ent{q_\Theta}
\qquad\text{evaluated at } q=(q_\Theta^\star,\{q_n\}),
\label{eq:logZ-identity}
\end{equation}
where the $q_n$ are the (fixed) residue factors used to build the potentials.
Consequently
\begin{equation}
\boxed{\;\mathcal L(q_\Theta^\star,\{q_n\}) \;=\; \log Z_{\mathrm{field}}
   \;+\; \sum_{n=1}^m \Ent{q_n}.\;}
\label{eq:elbo-eval}
\end{equation}
\end{proposition}
\begin{proof}
Write $E(\btheta,\bDelta)$ for the total energy, i.e.\ the sum of
\eqref{eq:field-root}--\eqref{eq:field-edge1} over the root and all edges. By
construction $E(\btheta,\bDelta)=\Eq{\prod_n q_n}{\log P(\btheta,\bDelta,\bx)}$
--- the potentials are exactly the residue-averaged log-joint as a function of
the field configuration. For a tree MRF, the Gibbs law
$q_\Theta^\star=e^{E}/Z_{\mathrm{field}}$ obeys
$\Ent{q_\Theta^\star}=-\Eq{q_\Theta^\star}{E}+\log Z_{\mathrm{field}}$, i.e.\
$\log Z_{\mathrm{field}}=\Eq{q_\Theta^\star}{E}+\Ent{q_\Theta^\star}$. Taking
$\Eq{q_\Theta^\star}{E}=\Eq{q_\Theta^\star}{\Eq{\prod_n q_n}{\log P}}
=\Eq{q}{\log P}$ gives \eqref{eq:logZ-identity}; adding
$\sum_n\Ent{q_n}$ and using \eqref{eq:entropy-split},\eqref{eq:elbo-def} gives
\eqref{eq:elbo-eval}. \texttt{field\_logbp} returns exactly the tree-MRF
$Z_{\mathrm{field}}$ and marginals, so the identity holds numerically.
\end{proof}

\subsection{Monotonicity and cost}

\begin{proposition}[Monotone coordinate ascent]
Let $q_\Theta^{t}=\arg\max_{q_\Theta}\mathcal L(q_\Theta,\{q_n^{t-1}\})$
(field update) and $q_n^{t}=\arg\max_{q_n}\mathcal L(q_\Theta^{t},\cdot)$
(residue update). Define the reported value
$R_t := \mathcal L(q_\Theta^{t},\{q_n^{t-1}\}) = \log Z_{\mathrm{field}}^{t}
 + \sum_n\Ent{q_n^{t-1}}$ (Prop.~\ref{prop:field-elbo}). Then
$R_t\le R_{t+1}\le\log P(\mathrm{obs})$ for all $t$.
\end{proposition}
\begin{proof}
Each block update is an exact maximiser, so
$\mathcal L(q_\Theta^{t},\{q_n^{t-1}\})\le\mathcal L(q_\Theta^{t},\{q_n^{t}\})
\le\mathcal L(q_\Theta^{t+1},\{q_n^{t}\})=R_{t+1}$, and $R_t\le\log P$ by
\eqref{eq:elbo-def}.
\end{proof}
Precomputing the fixed kernels ($\beta,W,\Psub$) once, each sweep costs
$O(N L^2)$ for the field BP plus $O(m N A^2)$ for the $m$ residue BPs and the
potential assembly: \textbf{linear in $m$}, never $A^m$. Contrast the exact
peel at $O(N L A^{2m})$.

\paragraph{Warm start.}
Initialise $q_\Theta$ to the field prior (potentials
\eqref{eq:field-root}--\eqref{eq:field-edge1} with the residue-evidence terms
removed) and run one residue update; this yields independent, field-prior-
averaged LG residue posteriors, from which coordinate ascent proceeds.

\section{Tightness and exactness}
\label{sec:tightness}

By \eqref{eq:elbo-def}, the gap is
$\log P(\mathrm{obs})-\mathcal L(q)=\KL{q}{P(\cdot\mid\mathrm{obs})}\ge0$, so
$\mathcal L$ is exact iff the true posterior lies in the family
\eqref{eq:family}: it must factor as [tree-Markov in $(\btheta,\bDelta)$]
$\times\prod_n$[tree-Markov in $\bx_n$], with the field independent of the
residues \emph{a posteriori}. The one approximation the family makes is
precisely to sever this field--residue seam.

\begin{proposition}[Exactness: field-independent emissions AND no jumps]
\label{prop:exact}
The product-of-trees bound is exact, $\mathcal L^\star=\log P(\mathrm{obs})$,
for any tree and any $m$ whenever both hold:
\begin{itemize}
\item[(a)] emissions are field-independent: $\pi(c_n,\theta,\cdot)$ and
      $Q(c_n,\theta)$ do not depend on $\theta$; and
\item[(b)] there are no jumps: $\beta_\tau\equiv1$, $W_\tau\equiv0$ --- i.e.\
      $\rho_{\mathrm{chain}}=0$, or the degenerate $L=1$ (single field state,
      $\rho$ a point mass), which forces (a) and (b) simultaneously.
\end{itemize}
\end{proposition}
\begin{proof}
Under (b) every edge has $\Delta_v=0$ and $\theta_v=\theta_r$ for all $v$: the
field is frozen at the root state. Under (a) the surviving branch factor
$\prod_n\Psub^{c_n}(\tau_v)(x_{n,u},x_{n,v})$ is field-free, so
\eqref{eq:joint} becomes
$\rho(\theta_r)\times\prod_n\big[\pi(c_n,x_{n,r})\prod_v\Psub^{c_n}(x_{n,u},x_{n,v})\big]$,
a product of a field prior and $m$ independent residue tree chains. Its
posterior has $\theta_r$ independent of the residues and each site an
independent $A$-state tree chain --- exactly the family \eqref{eq:family}, so
$\KL{q}{P(\cdot\mid\mathrm{obs})}=0$ at the block maximisers. Dropping (a) but
keeping (b) makes the frozen field a \emph{mixture}
$\sum_\theta\rho(\theta)\prod_n P(\bx_n\mid\theta)$ whose posterior correlates
$\theta_r$ with the residues, so (a) is genuinely required.
\end{proof}

\begin{remark}[The jump reset couples residues to $\bDelta$; $m=1$ is not exact]
Field-independent emissions alone (without (b)) do \emph{not} suffice: a jump
$\Delta_v=1$ \emph{resets} the residue by drawing it afresh from $\pi(c_n,\cdot)$
(a renewal), so the residue trajectory's per-edge transition depends on the
jump indicator $\Delta_v$ --- and hence on $q_\Theta$ --- even when the
emission distribution does not depend on the field \emph{value}. The
along-site coupling to $\bDelta$ is the seam the family severs. Consequently
the bound is inexact whenever there is jump uncertainty coupled to residue
data, \emph{including for a single site} $m=1$ at $\rho_{\mathrm{chain}}>0$.
This is the sharp contrast with the field-conditioned bounds
\texttt{elbo\_persite}/\texttt{elbo\_treestruct}, exact for $m=1$ (any depth)
and cherries (any $m$) because they retain the seam. In particular
$\rho_{\mathrm{chain}}=0$ with field-\emph{dependent} emissions is a strict
(loose) bound --- the frozen-field mixture above --- whereas the seam-keeping
bounds are exact there. Empirically the residual gap scales $\propto
\rho_{\mathrm{chain}}$ as $\rho_{\mathrm{chain}}\to0$ under (a).
The evaluation (\texttt{analysis/product\_of\_trees\_elbo\_eval.md}) reports
exactly how loose the bound is on real 128-leaf trees, and whether coordinate
ascent recovers enough of the gap to make it usable as a scorer.
\end{remark}

\begin{remark}[Relation to the previous bounds]
The three bounds form a ladder of families over the same joint:
node-mean-field (\texttt{elbo}) $\subset$ product-of-trees (this note)
partially-overlapping-with field-conditioned per-site
(\texttt{elbo\_persite}) $\subset$ full residue-correlated
(\texttt{elbo\_treestruct}). Product-of-trees restores the \emph{along-tree}
residue correlation that the node-mean-field floor crushes (every $q_n$ is a
full tree chain, not per-node marginals), while spending only $O(A^2)$ per
site; it drops the \emph{field--residue} correlation that
\texttt{elbo\_persite} keeps. Which of product-of-trees and
\texttt{elbo\_persite} is tighter is therefore case-dependent and is settled
empirically.
\end{remark}

\end{document}
